\documentclass[11pt]{article}
\usepackage[utf8]{inputenc}
\usepackage[T1]{fontenc}
\usepackage{geometry}
\usepackage{amsmath, amssymb}
\usepackage{setspace}
\geometry{a4paper, margin=2cm}

\begin{document}

\begin{center}
{\LARGE \textbf{Lista de Exercícios – P.A. e P.G. (Com Respostas)}}
\end{center}

\bigskip

\textbf{1.} Numa progressão aritmética (P.A.) cujo primeiro termo é $a_1=4$ e razão $r=6$, determine o 6.º termo. Mostre todos os passos.

\textbf{Resposta:}\\
O termo geral é $a_n = a_1 + (n-1)r$. \\
$a_6 = 4 + (6-1)\cdot 6 = 4 + 30 = 34$.\\
\textbf{Portanto, $a_6 = 34$.}

\bigskip

\textbf{2.} Obtenha a fórmula do termo geral $a_n$ da P.A. $(1,4,7,10,\ldots)$. Explique como chegou à expressão.

\textbf{Resposta:}\\
Identificamos $a_1=1$ e a razão $r = 3$. \\
Pela fórmula: $a_n = a_1 + (n-1)r = 1 + 3(n-1) = 3n - 2$. \\
\textbf{Logo, $a_n = 3n - 2$.}

\bigskip

\textbf{3.} Calcule a soma dos 12 primeiros termos da P.A. $(5,9,13,\ldots)$.

\textbf{Resposta:}\\
$a_1=5$, $r=4$. Primeiro: $a_{12} = 5 + (12-1)4 = 49$. \\
Soma: $S_{12} = \frac{12}{2}(5+49) = 6 \cdot 54 = 324$.\\
\textbf{Portanto, $S_{12}=324$.}

\bigskip

\textbf{4.} Que número deve ser adicionado à P.A. $(2,5,8,\ldots)$ para que se torne $(7,10,13,\ldots)$?

\textbf{Resposta:}\\
Comparando: $7-2=5$, $10-5=5$, $13-8=5$. \\
\textbf{Deve-se adicionar 5 a cada termo.}

\bigskip

\textbf{5.} Determine a razão da P.A. $(-10,-5,0,5,10,\ldots)$. Mostre o procedimento.

\textbf{Resposta:}\\
$r = -5 - (-10) = 5$. Confirmando: $0 - (-5)=5$. \\
\textbf{Logo, $r=5$.}

\bigskip

\textbf{6.} Numa P.G. com $a_1=3$ e razão $q=4$, encontre o 5.º termo.

\textbf{Resposta:}\\
$a_n = a_1 q^{n-1}$. \\
$a_5 = 3 \cdot 4^4 = 3 \cdot 256 = 768$. \\
\textbf{Portanto, $a_5 = 768$.}

\bigskip

\textbf{7.} Numa P.G. com $a_1=64$ e $q=\tfrac14$, calcule o 4.º termo.

\textbf{Resposta:}\\
$a_4 = 64 \left(\frac14\right)^3 = 64 \cdot \frac{1}{64} = 1$. \\
\textbf{Logo, $a_4 = 1$.}

\bigskip

\textbf{8.} Determine a razão da P.G. $(5,15,45,135,\ldots)$.

\textbf{Resposta:}\\
$q = \frac{15}{5} = 3$. Verificando: $\frac{45}{15}=3$. \\
\textbf{Assim, $q=3$.}

\bigskip

\textbf{9.} Calcule a soma dos 6 primeiros termos da P.G. $(1,2,4,8,\ldots)$.

\textbf{Resposta:}\\
$a_1=1$, $q=2$. \\
$S_6 = \frac{a_1(q^6 - 1)}{q - 1} = \frac{2^6 - 1}{1} = 63$. \\
\textbf{Logo, $S_6 = 63$.}

\bigskip

\textbf{10.} O produto dos três primeiros termos de uma P.G. é $27$ e a razão é $2$. Determine o primeiro termo.

\textbf{Resposta:}\\
Se $a_1 = x$, termos: $x,\;2x,\;4x$. Produto: $x \cdot 2x \cdot 4x = 8x^3 = 27$. \\
$x^3 = \frac{27}{8}$, logo $x = \frac{3}{2}$. \\
\textbf{Portanto, $a_1 = \frac{3}{2}$.}

\end{document}

